
\section{Related Work}
\subsection{Human Motion Datasets with Text Annotation}

The rapid progress of text-driven human motion modeling has been closely tied to the emergence of motion datasets with language annotation~\cite{bensabath2024cross,wu2025mg,wang2024scaling,guo2025snapmogen}. Early benchmarks such as KIT-ML~\cite{Plappert2016} and HumanML3D~\cite{guo2022generating} established the standard text-to-motion setting by providing paired motion clips and clip-level natural language descriptions, enabling shared embedding learning, retrieval-based training, and standardized evaluation~\cite{petrovich2023tmr}. 

Subsequent datasets broadened the supervision spectrum beyond coarse clip-level captions. BABEL~\cite{punnakkal2021babel} augments mocap sequences with sequence- and frame-level action labels, supporting temporally localized understanding. Motion-X~\cite{lin2023motion} provides expressive whole-body motion together with sequence-level semantics and video-derived frame-level pose descriptions. More recently, SnapMoGen~\cite{guo2025snapmogen} introduces long-form narrative descriptions with improved temporal continuity, while MotionLib~\cite{wang2024scaling} scales motion--text supervision to millions of samples with hierarchical annotations. Dense Motion Captioning further introduces CompMo, a large-scale dataset with temporally bounded multi-action sequences and grounded captions~\cite{xudense}. In parallel, FrankenMotion provides temporally aware part-level prompts for compositional motion modeling~\cite{li2026frankenmotion}.

Despite these advances, most motion--language datasets still provide supervision either at the clip level or through limited predefined labels, leaving many-to-many frame-level (or token-to-frame) correspondence under-explored. 

This gap motivates learning objectives that explicitly model fine-grained temporal correspondence between motion and language~\cite{xudense,li2025unimotion}, complementary to efforts that focus on stronger generators or clip-level supervision~\cite{tevet2022human,zhang2023generating,guo2024momask,meng2025rethinking}.

\subsection{Human Motion Understanding}
Human motion understanding has long been studied in computer vision through video-based tasks such as action recognition and temporal localization, where models align motion with action labels or textual descriptions~\cite{bertasius2021space,arnab2021vivit,nag2022proposal,feichtenhofer2022masked}. In contrast, our work focuses on structured human motion data (e.g., mocap sequences), where the temporal evolution of poses is directly available and can be paired with language supervision.

In the motion domain, early understanding models largely relied on shared embedding learning and retrieval-style objectives. MotionCLIP~\cite{tevet2022motionclip} aligns motion latents with the CLIP~\cite{radford2021learning} embedding space to exploit strong language priors, while TMR~\cite{petrovich2023tmr} combines contrastive learning with a motion synthesis objective for text--motion retrieval. Cross-dataset studies~\cite{bensabath2024cross} further show that such retrieval performance is sensitive to shifts in motion distributions and annotation styles, motivating robustness and transfer-aware evaluation.

More recent works move beyond global alignment toward fine-grained motion--language understanding. Dense Motion Captioning~\cite{xudense} requires temporally localizing and describing multiple actions within long motion sequences, while unified motion--language models such as MG-MotionLLM~\cite{wu2025mg} and UniMotion~\cite{li2025unimotion} further connect motion understanding, localization, and generation. At the same time, advances in generation and control, including ReMoDiffuse~\cite{zhang2023remodiffuse}, FineMoGen~\cite{zhang2023finemogen}, CoMo~\cite{huang2024controllable}, and recent text-conditioned generators~\cite{zhang2023generating,guo2024momask,meng2025rethinking,Zhao:DartControl:2025,barquero2024seamless,sun2024coma,jiang2023motiongpt,pinyoanuntapong2024mmm,dang2026segmo}, highlight the growing demand for fine-grained temporal correspondence between motion and language. Our work addresses this need by learning dense motion--text alignment from clip-level supervision.

\subsection{Semi-supervised Learning Using Optimal Transport}
Optimal transport (OT) provides a principled tool to establish soft correspondences between two sets of elements under marginal constraints, and has been widely used for matching and alignment problems~\cite{solomon2018optimal,peyre2019computational,cuturi2013sinkhorn}. Recent work has incorporated OT into representation learning and robust alignment.
For instance, OT-based matching has been used to mitigate data poisoning in CLIP-style~\cite{radford2021learning} pre-training by aligning samples through transport plans~\cite{zhang2025pre}.
OT has also been applied to understand and generalize CLIP by modeling relationships between sets of features via transport~\cite{shi2024ot, yao2021filip}. In this work, we focus on \emph{fine-grained correspondences} in the motion-language domain.